% SHAPE_THEOREM_SECTION.tex — reviewed version, integrated verbatim into 04-papers/46-row-d-dual-certificates/main.tex
% as \section{The shape of a proof} (sec:shape) before \section{Open problems}; Problem (R) prob:region added to sec:open.
% Corrections of the review SHAPE_THEOREM_REVIEW.md applied: three-tier theorem; O6 (Euler-preserving deletion) and the
% twist proposition proved in the paper; Beurling clause ARGUED; margin/universal criteria structural; constants T0(0)=0.74921, kappa_* >= 1.343.
% Compiles inside the preamble of main.tex; bibitems Paper47 and Titchmarsh1986 added there.

\section{The shape of a proof}\label{sec:shape}

This section collects the barriers of \S\ref{sec:negative}, \S\ref{sec:barriers} and \S\ref{sec:wh} into one statement about
the form that any valid proof of \eqref{eq:target} must have. Nothing here is a new route and nothing moves $\lamA$. Every
clause is either an instance of Lemma~\ref{lem:transfer} with an explicit falsifier, or a constraint on the precision or the
exactness of a proof; every clause excludes a kind of deduction, never the conclusion. Each statement carries one of the labels
\emph{proved}, \emph{cited}, \emph{conditional}, \emph{argued}, \emph{open} or \emph{numerical}; a clause that is only argued or
only numerical is stated as such and is not used in the theorem.

\subsection{Form objects, modifications and classes}

\begin{definition}[form objects]\label{def:shape-obj}
A \emph{comb} is a function $a:\N_{\ge2}\to\R_{\ge0}$, finite on every window, with weights $v_n(a)=2a(n)/\sqrt n$; its window
form $\mathfrak a^{(a)}_T$, its bottom $\lambda_{\mathrm A}(T;a)$ on $\PT$ and the statement $\mathrm{WP}_T(a)$ are those of
\S\ref{subsec:finite}, and $a=\Lambda$ gives $\lamA$ and \eqref{eq:target}. A \emph{form object} is a family $q=(q_T)_{T>0}$ of
closed semibounded forms, $q_T$ defined on $\PT\cap\mathcal D_T$ or on a closed subspace of $L^2(\IT)$, with bottoms
$\lambda(T;q)$. Every comb is a form object through $q_T=\mathfrak a^{(a)}_T|_{\PT}$, and $q_0:=(\mathfrak a_T|_{\PT})_T$ is the
object of $\zeta$. The zero-side objects of \S\ref{subsec:bar-fam} and \S\ref{subsec:bar-pat} are form objects that are not combs:
the quartet models $\mathfrak a_T+kQ_z$ of Definition~\ref{def:quartetmodel}, and the deformations $f_t$ of
Proposition~\ref{prop:ft}, whose weights have both signs.
\end{definition}

\begin{definition}[modifications]\label{def:shape-ops}
The following maps on form objects are used below.
\begin{enumerate}
\item[$\mathrm O_1$] finite deletion, $a\mapsto a\mathbf 1_{n\ge N}$;
\item[$\mathrm O_2$] one-term amplification, $a\mapsto a+c\mathbf 1_{n=n_0}$;
\item[$\mathrm O_3$] tail deletion, $a\mapsto a\mathbf 1_{n<n_0}$;
\item[$\mathrm O_4$] uniform rescaling, $a\mapsto ca$ with $c>0$, and its small form $\mathrm O_4^\varepsilon$ with $c\in(1-\varepsilon,1+\varepsilon)$;
\item[$\mathrm O_5$] a change of one weight by $\varepsilon$;
\item[$\mathrm O_6$] Euler-preserving prime deletion, $a\mapsto a_P:=\Lambda\mathbf 1\{n\text{ is a power of a prime }p\ge P\}$;
\item[$\mathrm O_7$] multiplication by quartets, $\xi\mapsto\xi_{z,k}$ (Definition~\ref{def:quartetmodel});
\item[$\mathrm O_8$] self-dual deformation, $\zeta\mapsto f_t=\zeta+tg$ (Lemma~\ref{lem:selfdual});
\item[$\mathrm O_9$] a functional-equation-preserving twist $\xi\mapsto\xi R$ by a rational function (Proposition~\ref{prop:shape-twist});
\item[$\mathrm O_{10}$] Euler removal: passage to an Epstein zeta function of class number greater than one (Remark~\ref{rem:twohalves});
\item[$\mathrm O_{11}$] Beurling repositioning: $\N$ replaced by a generalized-integer system with a positive prime measure;
\item[$\mathrm O_{12}$] field change, $\xi\mapsto\xi_K$ for a number field $K$;
\item[$\mathrm O_{13}$] restriction of the family $T\mapsto\lambda(T;q)$ to a bounded set of windows $J\subset(0,T_1]$;
\item[$\mathrm O_{14}$] symbol truncation, $\sT\mapsto\sigma^{(\eta)}=\min(\sT,\eta)$ (Theorem~\ref{thm:nolocal});
\item[$\mathrm O_{15}$] Poisson randomization of the positions $\{\log n\}$ at the intensity of the primes (Theorem~\ref{thm:poisson}).
\end{enumerate}
Two further modifications produce no falsifier and enter only the structural part: the reflection $\sigma\mapsto1-\sigma$ of the
family $D_\sigma$ of Remark~\ref{rem:fe}, and a bounded perturbation $q_T\mapsto q_T+\mathfrak b$ with $\|\mathfrak b\|\le\lamA$
(Lemmas~\ref{lem:wh-T}, \ref{lem:wh-Tprime}).
\end{definition}

\begin{definition}[deduction classes]\label{def:shape-class}
A \emph{deduction class} $\mathcal A$ is specified by a set $\mathsf P$ of properties of form objects; a deduction in $\mathcal A$ is a
valid chain ``$q$ has $\mathsf P$, hence $\lambda(T;q)\ge0$'', for every $T$ or at a fixed $T$. The class is \emph{closed under} a
modification $\mathrm O$, for the parameters named in the corresponding lemma, if $q$ having $\mathsf P$ implies that $\mathrm O(q)$
has $\mathsf P$. The classes already named in this paper are: identities linear in the autocorrelation with Bochner or Cauchy--Schwarz
positivity (Proposition~\ref{prop:opcalc}); norm, moment, determinant, local, family, pattern and coherence functionals
(\S\ref{sec:wh}, \S\ref{sec:barriers}, Proposition~\ref{prop:wit-sens}); the zero-information class $\mathcal C_{\mathrm{zeros}}$ and
the classes $\mathcal C_{\mathrm{PNT}}(\varepsilon)$, $\mathcal C_{\mathrm{BT}}$ of Corollary~\ref{cor:void}; the progression classes of
Proposition~\ref{prop:prog}; the dual, lattice, flat-top and square certificates (Definition~\ref{def:PD}, Theorems~\ref{thm:Z},
\ref{thm:sos}); the symmetrizations of $D_\sigma$ (Remark~\ref{rem:fe}); and the axioms A1--A6 and the field-uniform properties of
\cite{Paper47}.
\end{definition}

\subsection{Transfer lemmas}

Each lemma names the modification, the class it closes, the falsifier, the windows at which the falsifier is negative, and the label.
The first group is proved in this paper.

\begin{lemma}[$\mathrm O_1$, finite deletion; proved]\label{lem:shape-del}
If $\mathsf P$ is preserved by $a\mapsto a\mathbf 1_{n\ge N}$, as it is for $\mathcal C_{\mathrm{PNT}}(\varepsilon)$,
$\mathcal C_{\mathrm{zeros}}$ and $\mathcal C_{\mathrm{BT}}$, then $\Lambda\mathbf 1_{n\ge N}$ with $N\ge\max(5,e^{2T})$ has $\mathsf P$
and fails $\mathrm{WP}_T$ for every $T\in[0.60,\frac12\log N]$, and for every $T>T_0(0)=0.74921$ (Theorem~\ref{thm:finmod}(i),
Corollary~\ref{cor:void}, Proposition~\ref{prop:arch}(ii)).
\end{lemma}

\begin{lemma}[$\mathrm O_2$, amplification; proved]\label{lem:shape-amp}
If $\mathsf P$ is preserved by $a\mapsto a+c\mathbf 1_{n=2}$, then $\Lambda+c\mathbf 1_{n=2}$ with $c>c_*(T,2)$ has $\mathsf P$ and fails
$\mathrm{WP}_T$ at every $T>\frac12\log2$, including the windows $(\frac12\log2,\frac12\log5]$ at which $\mathrm{WP}_T(\Lambda)$ holds
(Theorem~\ref{thm:finmod}(ii)).
\end{lemma}

\begin{lemma}[$\mathrm O_3$, tail deletion; proved]\label{lem:shape-tail}
If $\mathsf P$ holds for every $a$ with $a=\Lambda$ on $n<n_0$ and $0\le a\le\Lambda$ (the exact small values together with every upper-bound
sieve inequality), then $\Lambda\mathbf 1_{n<n_0}$ has $\mathsf P$ and fails $\mathrm{WP}_T$ for all large $T$; the dilated tests of
Proposition~\ref{prop:finsupp} have limit value $C_0-\sum_{n<n_0}v_n<0$ (Corollary~\ref{cor:wtail}(i)). The first failing window is not located.
\end{lemma}

\begin{lemma}[$\mathrm O_4$, rescaling; proved for $c$ large and small, conditional for $\mathrm O_4^\varepsilon$]\label{lem:shape-scale}
If $\mathsf P$ is preserved by $a\mapsto ca$ for every $c>0$ (an Euler product $\zeta^c$, the support, the relative sizes of the weights),
then $c\Lambda$ fails $\mathrm{WP}_T$ at every $T>\frac12\log2$ for $c$ large, and at every $T\ge0.60$ for $c$ small. If $\mathsf P$ is
preserved only by $c\in(1-\varepsilon,1+\varepsilon)$, then $(1+\operatorname{sgn}(\kappa_T)\varepsilon)\Lambda$ fails $\mathrm{WP}_T$ as soon as
$\varepsilon>\lamA/|\kappa_T|$, \emph{provided} $\kappa_T\ne0$, where $\kappa_T=\sum_nv_nc_h(\log n)$ for a real unit minimizer $h$
(Proposition~\ref{prop:wit-sens}); whether $\kappa_T\ne0$ at some window is not known, and $\kappa_{\log2}\approx0.09$ is a numerical value
(Remark~\ref{rem:wit-sens}).
\end{lemma}
\begin{proof}
For $c$ large take $F_+=\phi+\phi(\cdot-\log2)$ of Lemma~\ref{lem:cpos} with $u_0=\log2$ and $\ell<\log\frac32$. The autocorrelation
$c_{F_+}$ is supported in $\{|u|<\ell\}\cup\{|u\mp\log2|<\ell\}$, which contains no $\log n$ with $n\ge3$, and $c_{F_+}(\log2)=\|\phi\|^2>0$.
Hence $\mathfrak a^{(c\Lambda)}_T(F_+)=\mT(F_+)-cv_2\|\phi\|^2\to-\infty$ as $c\to\infty$. For $c$ small, the function $\mu_T$ of
Lemma~\ref{lem:degen} is finite and concave on $\R^{\NT}$, hence continuous, and $\mu_T(0)=\laminfP(T)<0$ for $T\ge0.60$
(Corollary~\ref{cor:Tarch}); so $\lambda_{\mathrm A}(T;c\Lambda)=\mu_T(cv)<0$ for $c$ small. The small form is Proposition~\ref{prop:wit-sens}.
\end{proof}

\begin{lemma}[$\mathrm O_5$, one weight; conditional]\label{lem:shape-weight}
Let $F_\star$ be a real unit minimizer at $T$ and $n_0\in\NT$ with $c_{F_\star}(\log n_0)\ne0$. If $\mathsf P$ is preserved by changing the
weight at $n_0$ by $\varepsilon$, then $\Lambda+s\varepsilon\mathbf 1_{n=n_0}$ with $s=\operatorname{sgn}c_{F_\star}(\log n_0)$ fails $\mathrm{WP}_T$
for every admissible $\varepsilon>\varepsilon_*(T,n_0)=\sqrt{n_0}\,\lamA/(2|c_{F_\star}(\log n_0)|)$, where for $s=-1$ admissibility means
$\varepsilon\le\Lambda(n_0)$ (Remark~\ref{rem:resweights}). If every real minimizer has $c_{F_\star}=0$ on $\log\NT$, the clause is empty
and Lemma~\ref{lem:degen} describes that case; whether this happens at any window is not known. The numerical scale $10^{-5}$ at $T=\log2$
(Remark~\ref{rem:wit-sens}) is not certified.
\end{lemma}

\begin{proposition}[$\mathrm O_6$, Euler-preserving prime deletion; proved]\label{prop:shape-eulerdel}
Let $P\ge3$ and $a_P:=\Lambda\mathbf 1\{n\text{ is a power of a prime }p\ge P\}$. Then $a_P\ge0$ has the generating function
$\zeta_{a_P}(s)=\prod_{p\ge P}(1-p^{-s})^{-1}=\zeta(s)\prod_{p<P}(1-p^{-s})$, an Euler product over $\N$ with nonnegative integral
Dirichlet coefficients; $\zeta_{a_P}$ has exactly the zeros of $\zeta$ in $0<\Rel s<1$ and the pole at $s=1$; $a_P(p^k)=\log p$ on its
support; $a_P\le\Lambda$; and $a_P$ fails $\mathrm{WP}_T$ for every $T\in[0.60,\frac12\log P]$. Of the properties listed in
Definition~\ref{def:shape-class}, $a_P$ breaks only the functional equation.
\end{proposition}
\begin{proof}
$\sum_na_P(n)n^{-s}/\log n=\sum_{p\ge P}\sum_k p^{-ks}/k=-\sum_{p\ge P}\log(1-p^{-s})$, which gives the Euler product; the coefficients of
$\prod_{p\ge P}(1-p^{-s})^{-1}$ count factorizations, so they are nonnegative integers. The finite product $\prod_{p<P}(1-p^{-s})$ is
entire and vanishes only on $\Rel s=0$. The comb of $a_P$ is empty on $\NT$ when $e^{2T}\le P$, so $\mathrm{WP}_T(a_P)$ is
$\laminfP(T)\ge0$, false for $T\ge0.60$ (Corollary~\ref{cor:Tarch}). Finally $\zeta_{a_P}(1-s)\ne\zeta_{a_P}(s)\,\Gamma$-factor, since
$\prod_{p<P}(1-p^{-s})$ is not invariant under $s\mapsto1-s$.
\end{proof}

\begin{lemma}[$\mathrm O_7$, quartets; proved]\label{lem:shape-quartet}
The models $\xi_{z,k}$ keep the comb $(\Lambda(n))_n$ exactly, the functional equation, reality and the pole, and share with $\zeta$ the
properties of Proposition~\ref{prop:quartetmodel}; for $k\ge k_0(T,\delta,\gamma)$ they have a witness at $T$, hence at every $T'\ge T$
(Theorem~\ref{thm:wit-model}(iii)). Their arithmetic side differs from that of $\zeta$ only by the absolutely continuous part $-k\beta_z$.
\end{lemma}

\begin{lemma}[$\mathrm O_8$, self-dual deformation; proved]\label{lem:shape-ft}
$f_t=\zeta+tg$ keeps every $\Lambda(n)$ with $n<\lambda_1$ and Riemann's functional equation with the archimedean factor and pole of
$\zeta$, and for $|t|$ large has a primitive witness at the window $\frac12\log\lambda_1+\varepsilon$ (Propositions~\ref{prop:ft},
\ref{prop:ftprim}). Its frequencies $m/\sqrt M$ are not integers: it breaks the Euler product over $\N$.
\end{lemma}

\begin{proposition}[$\mathrm O_9$, a twist that keeps the line data; proved]\label{prop:shape-twist}
Fix $\gamma_0>0$ with $\zeta(\frac12+i\gamma_0)\ne0$ and $0<\delta<\frac18$, and put
\[
R(s)=\prod_{\pm}\frac{(s-\frac12-\delta\mp i\gamma_0)(s-\frac12+\delta\mp i\gamma_0)}{(s-\frac12-2\delta\mp i\gamma_0)(s-\frac12+2\delta\mp i\gamma_0)},
\qquad\tilde\xi:=\xi R,\quad\tilde\zeta:=\zeta R,
\]
and $\ell(\tau):=\log|\xi(\frac12+i\tau)|$, $\tilde\ell(\tau):=\log|\tilde\xi(\frac12+i\tau)|$. Then:
\emph{(i)} $\tilde\xi(1-s)=\tilde\xi(s)$ and $\tilde\xi$ is real on $\R$; $\tilde\xi$ is meromorphic with four poles, at
$\frac12\pm2\delta\pm i\gamma_0$, inside the strip.
\emph{(ii)} On the line, $R(\frac12+i\tau)=\prod_\pm\frac{(\tau\mp\gamma_0)^2+\delta^2}{(\tau\mp\gamma_0)^2+4\delta^2}\in[\frac1{16},1]$; hence
$|\tilde\zeta|\le|\zeta|$ on the line, and, off zeros and poles, $\Rel\frac{\tilde\zeta'}{\tilde\zeta}(\frac12+i\tau)=-\frac12m(\tau)$,
the same identity as for $\zeta$.
\emph{(iii)} $\tilde\ell=\ell+\log R(\frac12+i\,\cdot)$ with $\int_\R|\log R(\tfrac12+i\tau)|\,d\tau=4\pi\delta$.
\emph{(iv)} $\tilde\ell''(\gamma_0)=\ell''(\gamma_0)+\frac3{2\delta^2}+O(\delta^2\gamma_0^{-4})$.
Consequently, let $\mathcal M$ be the class of arguments whose inputs about $\zeta$ on the line are: the functional equation with
$\Gamma_\R$ and conductor $1$, reality, the identity in (ii), and finitely many inequalities of the forms
$\int_I\varphi\,\ell\,d\tau<B$ or $\int_I|\zeta(\frac12+i\tau)|^{2k}d\tau<B_k$ with fixed bounded $I$, bounded $\varphi$ and strict
inequality. No argument in $\mathcal M$ bounds $\ell''$ from above on an interval containing $\gamma_0$. The class does not contain
arguments that use the holomorphy of $\xi$ in the strip together with its order-one growth.
\end{proposition}
\begin{proof}
(i) The zero and pole sets of $R$ are invariant under $s\mapsto1-s$ and $s\mapsto\bar s$, and a direct check of the eight linear factors
gives $R(1-s)=R(s)$ and $\overline{R(\bar s)}=R(s)$; $\xi$ has both symmetries. (ii) On $s=\frac12+i\tau$ each factor has modulus
$((\tau\mp\gamma_0)^2+\delta^2)/((\tau\mp\gamma_0)^2+4\delta^2)\in[\frac14,1]$ and the product is real positive. A function real on the
line has purely imaginary logarithmic derivative there, so $\Rel\frac{\tilde\xi'}{\tilde\xi}(\frac12+i\tau)=0$; since
$\Rel[\frac1s+\frac1{s-1}]=0$ on the line and $\Rel[\frac12\psi(\frac s2)-\frac12\log\pi]=\frac12m(\tau)$ by \eqref{eq:m}, this is the
identity stated, for $\tilde\zeta$ exactly as for $\zeta$. (iii) $\int_\R[\log(x^2+a^2)-\log(x^2+b^2)]dx=2\pi(a-b)$ for $a,b>0$, applied
with $(a,b)=(\delta,2\delta)$ to each of the two factors; $\log R\le0$ on the line. (iv) $\frac{d^2}{dx^2}\log(x^2+a^2)|_{x=0}=2/a^2$, so
the factor centred at $\gamma_0$ contributes $2/\delta^2-2/(4\delta^2)=3/(2\delta^2)$ to the second derivative at $\tau=\gamma_0$, and the
factor centred at $-\gamma_0$, evaluated at distance $2\gamma_0$, contributes $O(\delta^2\gamma_0^{-4})$. For the consequence: suppose an
argument in $\mathcal M$ derived $\ell''\le K$ on an interval $J\ni\gamma_0$. By (i), (ii) the first three inputs hold for $\tilde\zeta$ for
every $\delta$; by (ii) and (iii) each inequality of the last kind holds for $\tilde\zeta$ once $\delta$ is small (the moments of
$|\tilde\zeta|^{2k}$ do not exceed those of $|\zeta|^{2k}$, and $\int_I\varphi\tilde\ell$ differs from $\int_I\varphi\ell$ by at most
$4\pi\delta\|\varphi\|_\infty$). So the argument gives $\tilde\ell''\le K$ on $J$, contradicting (iv) for $\delta$ small. The last
sentence is (i): $\tilde\xi$ is not entire.
\end{proof}

No negative value of a window form of $\tilde\zeta$ on $\PT$ has been computed; Proposition~\ref{prop:shape-twist} excludes a class of
pointwise arguments about the curvature of $\ell$ on the line, and is not a falsifier in the sense of Lemma~\ref{lem:transfer}. An
$m$-fold twist with a witness at every window, in the manner of Theorem~\ref{thm:wit-model}(iii), would make it one; the form
$Q_{\delta+i\gamma_0}-Q_{2\delta+i\gamma_0}$ of Lemma~\ref{lem:quartet} is the natural candidate and has not been examined.

The second group of transfer lemmas rests on inputs that are cited, argued or open.

\begin{lemma}[$\mathrm O_{10}$, Euler removal; cited for the complete form, argued for the primitive bottom]\label{lem:shape-epstein}
Let $Q$ be a positive definite binary quadratic form whose discriminant has class number greater than one, $Z_Q$ its Epstein zeta
function and $\Xi_Q$ its completion. $\Xi_Q$ shares with $\xi$ the functional equation with $\Gamma(s/2)$, reality, the polar set and
nonnegative Dirichlet coefficients, that is, the axioms A1--A6 of \cite{Paper47}, and has zeros in $\Rel s>1$
\cite{DavenportHeilbronn1936}. By Weil's criterion applied to $Z_Q$ \cite{Weil52}, the Weil form of $Z_Q$ takes a negative value on a
test of compact support; hence the complete window form built from the coefficients of $-Z_Q'/Z_Q$ (which have both signs: there is no
Euler product) and the archimedean factor of $\Xi_Q$ is negative at some window and at all larger ones. A witness on the primitive
subspace $\PT$, as in Proposition~\ref{prop:ftprim} for $f_t$, has not been written. (Remark~\ref{rem:twohalves}; the falsifier and the
axiom list are those of \cite[Thms.~18.10, 18.16]{Paper47}, which concern a different inequality and are used here only as a source.)
\end{lemma}

\begin{lemma}[$\mathrm O_{11}$, Beurling systems; argued]\label{lem:shape-beurling}
There are generalized-integer systems with nonnegative prime measure, an Euler product and a counting function $N(x)=\rho x+O(x^\theta)$
whose zeta function has zeros off the line $\Rel s=\frac12$ \cite{DiamondZhang2016}. Such a system keeps every property of $\Lambda$ that is
stated inside the semigroup (nonnegativity, the Euler product, the identity $\Lambda*1=\log$, the sieve axioms) and breaks only the
integrality of the frequencies. It has no functional equation with $\Gamma_\R$, and no window form of it has been shown negative: the
step from an off-line zero to a negative window bottom would require an explicit formula for the system with a test of compact support,
which is not written. The lemma is therefore an argued clause and is not used in Theorem~\ref{thm:shape}(A).
\end{lemma}

\begin{lemma}[$\mathrm O_{12}$, field change; transfer principle proved, falsifier open]\label{lem:shape-field}
A derivation of $\lamA\ge0$ that invokes only properties shared by the completed zeta functions $\xi_K$ of all number fields $K$ is a
derivation of the corresponding conclusion for every $\xi_K$. No window $T$ and field $K$ with a negative primitive bottom for $\xi_K$
is exhibited in this paper; the clause therefore says only that a proof must name the property that separates $\zeta$ from every
$\xi_K$, or else be held to the standard of proving the sign for all of them (the principle of \cite[Thm.~18.1]{Paper47}, whose own
proof passes through a statement about $\zeta_K$ that is not used here).
\end{lemma}

\begin{lemma}[$\mathrm O_{13}$, bounded families; proved]\label{lem:shape-fam}
If $\mathsf P$ depends on $T\mapsto\lambda(T;q)$ only through $J\subset(0,T_1]$, monotonicity and continuity, and on the weights only
through membership in a class closed under $\mathrm O_2$, then $\Lambda+c\mathbf 1_{n=n_0}$ with $n_0>e^{2T_1}$ has $\mathsf P$ and fails
$\mathrm{WP}_T$ for $T>\frac12\log n_0$ and $c>c_*(T,n_0)$; if instead the weights are only known to come from a self-dual Dirichlet
series with the functional equation of $\zeta$, $f_t$ does the same (Theorems~\ref{thm:famblind}, \ref{thm:famFE}). Not excluded: the use
of $\lambda$ on an unbounded set of windows, which enters only through $\lim_T\lamA\ge0$, and two-sided laws $\lamA=L(T)(1+o(1))$
with $L>0$ (Remark~\ref{rem:bar-fam}).
\end{lemma}

\begin{lemma}[$\mathrm O_{14}$, local criteria; proved]\label{lem:shape-local}
No criterion that depends on $\sT$ only through its restriction to $\{\sT<\eta\}$ decides the sign at any window: $\sigma^{(\eta)}$ has the
same negative set and a negative form for $\eta<\eta_0(T)$ (Theorem~\ref{thm:nolocal}).
\end{lemma}

\begin{lemma}[$\mathrm O_{15}$, random positions; proved]\label{lem:shape-poisson}
For $T\ge T(c,\kappa)$ the Poisson comb of intensity $\kappa$ fails $\mathrm{WP}_T$ with probability at least $1-e^{-cT\log T}$; in
particular some realization fails it, so no deduction that uses the positions $\{\log n\}$ only through their intensity and moments proves
$\mathrm{WP}_T$ at such $T$ (Theorem~\ref{thm:poisson}).
\end{lemma}

\begin{lemma}[derived classes; proved]\label{lem:shape-classes}
\emph{(i)} Classes defined by statements on $\psi_a(x;q,c)$ that are stable under $O(1)$ changes are closed under $\mathrm O_1$ and
$\mathrm O_2$ (Proposition~\ref{prop:prog}); classes $0\le a\le\Lambda$ under $\mathrm O_1$. \emph{(ii)} $\mathcal C_{\mathrm{zeros}}$ is
closed under $\mathrm O_1$ and $\mathrm O_2$, since $\zeta_a=\zeta E$ with $E$ entire and zero-free (Corollary~\ref{cor:void}); and no finite
symmetric multiset of ordinates carries a flat-top certificate (Theorem~\ref{thm:finZ}). \emph{(iii)} Deductions from window observables
(the bottom up to a relative error, the parity of the minimizer, the peak of $|\widehat{h^*}|^2$, eigenvalue counts, laws fitted to these)
are excluded numerically by $\xi_{z,1}$ and, for the shared properties of Proposition~\ref{prop:quartetmodel}, by $\xi_{z,k}$ with $k$ large
(Corollary~\ref{cor:pattern}).
\end{lemma}

\subsection{Structural lemmas}

These constrain every proof, in whatever class, and produce no falsifier.

\begin{lemma}[margin and universal criteria; proved, inputs cited]\label{lem:shape-margin}
\emph{(i)} $\lamW\le\lamP\le\lamA\le4\cdot10^9Te^{-15.449T}$ for $T\ge1$ (Corollary~\ref{cor:margin1}; the zero data are those of
\cite{PlattTrudgian2021}). Any lower bound that proves \eqref{eq:target} at $T$ has this absolute precision; relative to the comb mass
$\Aconst\asymp e^T$ it is $e^{-16.4T}$. \emph{(ii)} A criterion valid for every nonnegative trace-class operator, or for every bounded
perturbation of the form, has margin at most $k\lamA/c_P$ (moments), $L\lamA$ (Lipschitz functionals), and a leading term of the window
determinant with relative error $\varepsilon_0\ge1$ never decides (Lemmas~\ref{lem:wh-T}, \ref{lem:wh-Tprime},
Theorem~\ref{thm:wh-detinherit}(iii)). The modified objects in (ii) have bottom $0$, not a negative bottom: these are precision
constraints, not instances of Lemma~\ref{lem:transfer}.
\end{lemma}

\begin{proposition}[criticality]\label{prop:shape-crit}
\emph{(a) The centre (proved).} $m(0)=C_0=-2\zeta'(\frac12)/\zeta(\frac12)$ (Proposition~\ref{prop:boundary}), and the zero sum behind it
vanishes termwise: for a pair $\rho$, $1-\bar\rho$ one has $\Rel[\frac1{\frac12-\rho}+\frac1{\frac12-(1-\bar\rho)}]=0$. The identity holds
for every zero configuration symmetric under $\rho\mapsto1-\bar\rho$ and carries no slack.
\emph{(b) The packets (conditional).} For every $\gamma>0$ the threshold of $(E_\gamma)$ is exactly $m^\sharp(\gamma)$
(Proposition~\ref{prop:Egamma-exact}). If, in addition, the Riesz means of the twisted sums $\mathcal F_T(F_T)/\|F_T\|^2$ converge
whenever they are bounded, then $(E_\gamma)$ with any additive constant $C\ge0$ is equivalent to $(E_\gamma)$ itself: there is no slack at
any packet frequency. The additional hypothesis is a statement about the divergent alternative that is not proved here.
\emph{(c) The integer resolution (proved; a Mertens-type estimate cited).} Let $P_T(\tau)=\sum_{n\in\NT}v_n\cos(\tau\log n)$. Then
\[
\lim_{Y\to\infty}\frac1{2Y}\int_{-Y}^{Y}P_T(\tau)^2d\tau=\frac12\sum_{n\in\NT}v_n^2=4T^2\,(1+o(1))\qquad(T\to\infty),
\]
so the root mean square of the comb is $2T(1+o(1))$, while the archimedean symbol at the frequency $e^{2T}$ of the integer resolution is
$m(e^{2T})=2T+O(1)$ (Lemma~\ref{lem:mbounds}(ii)): comb and symbol are of the same size at that resolution.
\emph{(d) The density ceiling (proved).} A comb $a\ge0$ with $\lambda_{\mathrm A}(T;a)\ge0$ for every $T$ has
$\liminf S(x)/\log x\le1$, and $\Lambda$ sits at the ceiling (Theorem~\ref{thm:ceiling}); the main-term comb $\rho_\kappa$ is positive at
every window if and only if $\kappa\ge\kappa_\star$, where $\kappa_\star\ge\frac14|C_0|=1.343>1$ (Theorem~\ref{thm:kappa}). So the main
term of $\Lambda$ is below the positivity threshold of main terms.
\emph{(e) No asymptotic slack (proved).} \eqref{eq:target} holds as soon as $\lamA\ge-\varepsilon(T)$ for all $T$ with
$\varepsilon(T)\to0$, or $\lamA\ge0$ along one sequence $T_k\to\infty$ (Theorem~\ref{thm:wit-Tc}(iv)); there is no weaker asymptotic
sub-statement in $T$.
\end{proposition}
\begin{proof}
(a) Proposition~\ref{prop:boundary} and the displayed computation, in which $(\frac12-\beta)+(\beta-\frac12)=0$ for $\rho=\beta+i\gamma$.
(b) Proposition~\ref{prop:Egamma-exact}: under the hypothesis, a bounded $\limsup$ is a limit, and a limit that is at most
$m^\sharp(\gamma)+C$ for the packets $F_T$ at every scale is at most $m^\sharp(\gamma)$, since the limit does not depend on $C$.
(c) The frequencies $\log n$, $n\in\NT$, are distinct, so the Ces\`aro mean of $P_T^2$ is $\frac12\sum v_n^2=2\sum_{n<e^{2T}}\Lambda(n)^2/n$.
Since $\Lambda(n)^2=\Lambda(n)\log n$ on primes and $\sum_{p^k,k\ge2}\Lambda(p^k)\log p^k/p^k<\infty$, one has
$\sum_{n\le x}\Lambda(n)^2/n=\sum_{n\le x}\Lambda(n)\log n/n+O(1)=\frac12\log^2x+O(\log x)$ \cite[Thm.~2.7 and partial summation]{MontgomeryVaughan2007},
and $x=e^{2T}$ gives $2\cdot\frac12(2T)^2=4T^2$. (d) and (e) are the cited theorems.
\end{proof}

\begin{lemma}[antisymmetry; proved]\label{lem:shape-antisym}
$D_{1-\sigma}=-D_\sigma$ exactly, every combination of the family $D_\sigma$ symmetric under $\sigma\mapsto1-\sigma$ vanishes, and
$D_{1/2}\equiv0$ off the ordinates (Remark~\ref{rem:fe}). No positivity is obtained by symmetrizing the family.
\end{lemma}

\begin{lemma}[operator calculus and squares; proved]\label{lem:shape-lin}
\emph{(i)} Every expression $\sum_i\Rel\langle A_i*F,B_i*F\rangle$ with finite measures $A_i,B_i$ equals $\int_{(-2T,2T)}\Rel\cF\,d\nu$
(Proposition~\ref{prop:opcalc}): an identity in shift-multiplier calculus is an identity of measures on the window and carries the whole
margin, and the only further tool of that calculus is the positivity $\|A*F\|^2\ge0$. The statement covers convolution operators only.
\emph{(ii)} No arithmetic square decomposition exists at any $T\ge\frac12\log3$; with a free positive definite part the decompositions are
exactly the certificates of Definition~\ref{def:PD}, whose value is $\lamP$ (Theorems~\ref{thm:sos}, \ref{thm:D1}). The interval
$(\frac12\log2,\frac12\log3)$ is not decided.
\end{lemma}

\begin{lemma}[orders, relative bounds, exterior measures, termwise bounds; proved]\label{lem:shape-misc}
The M\"obius resolvent is not monotone for the order of $\mathcal K^*_T$ (Proposition~\ref{prop:cone}); no relative bound with constants
$(\theta,\kappa,C)$ decides for large $T$, not even $\theta=0$ (Proposition~\ref{prop:relative}); no certificate of the complete form with
a positive exterior measure exists at the windows where $\kappa_T(\Phi)<0$ (Theorem~\ref{thm:posext}; the two values in
Remark~\ref{rem:posext} are floating evaluations); and a bound that treats the comb term by term costs its full mass
(Proposition~\ref{prop:primeblind}).
\end{lemma}

\begin{lemma}[flows and new contacts; identity proved, readings argued]\label{lem:shape-flow}
The coefficient of any $C^1$ flow of $\log\det$ on Galerkin sections near a window with small bottom is dominated by
$\langle A'h^*,h^*\rangle/\lamA$, a functional of the minimizer (Lemma~\ref{lem:flowdet}, Remark~\ref{rem:bar-loc}); positivity at $T>T_0$
imposes on the new contacts an averaged inequality on $U_{T_0}$, never a pointwise one (Theorem~\ref{thm:X}), and the converse of that
inequality fails for the quartet model, numerically (Remark~\ref{rem:X}).
\end{lemma}

\begin{proposition}[the three inputs enter jointly]\label{prop:shape-inputs}
Consider the three properties: (E) the Euler product over $\N$, equivalently $\sum_{d\mid n}a(d)=\log n$ for $n\in\NT$; (F) the functional
equation with $\Gamma(s/2)$, conductor $1$, reality and the pole; (I) the integrality of the frequencies, that is, support on $\log\N$.
\emph{(i) (proved)} A deduction that uses (E) and (I) with nonnegativity, the local values $a(p^k)=\log p$ and the location of the zeros
and the pole, but not (F), proves nothing: $a_P$ has all of these (Proposition~\ref{prop:shape-eulerdel}).
\emph{(ii) (cited)} A deduction that uses (F) and (I) with nonnegative coefficients, but not (E), proves nothing for the complete form:
$\Xi_Q$ has all of these (Lemma~\ref{lem:shape-epstein}).
\emph{(iii) (argued)} A deduction that uses (E) and (F)-free semigroup properties, but not (I), is matched by the Beurling systems
of Lemma~\ref{lem:shape-beurling}.
\emph{(iv) (proved)} The modifications $\mathrm O_7$ and $\mathrm O_8$ keep (F): the quartet models keep (E) as a comb identity but add an
absolutely continuous arithmetic part, and $f_t$ breaks (I). Hence a property set that survives (i), (ii) and (iv) contains (E) and (F)
jointly, with (I); and (E) with (F) at conductor $1$ determines $\zeta$ up to a constant, by Hamburger's theorem \cite[\S2.13]{Titchmarsh1986}.
\end{proposition}
\begin{proof}
(i), (ii), (iv) are the lemmas cited; the quartet models have logarithmic derivative $\xi'/\xi$ plus a rational function, which is not a
Dirichlet series over $\N$, and the frequencies $m/\sqrt M$ of $f_t$ are irrational. The last sentence is Hamburger's theorem: a Dirichlet
series over $\N$ of finite order with the functional equation of $\zeta$ is a constant multiple of $\zeta$.
\end{proof}

\begin{lemma}[no weaker substatement; proved]\label{lem:shape-restate}
Each of the following is an exact restatement of \eqref{eq:target}, not a weaker statement: a flat-top certificate $f_w\ge0$ at every
window (Theorem~\ref{thm:Z}: $f_w=\sum_\rho\wh(\tau-\gamma_\rho)$, so its content is the explicit formula on the tests $w\cos(\tau\,\cdot)$);
the indeterminacy of the extension problem at every window off $\{\lamW=0\}$ (Theorem~\ref{thm:KEPdich}); the positive definiteness of
the one-profile function $\mathsf B$ (Theorem~\ref{thm:oneprofile}); the positivity of $\mathrm{Arch}_M-\mathrm{Div}_M-\mathrm{Near}_M$ for
every $M$ (Proposition~\ref{prop:AM}); the finite matrices $S_{T,k(T)}\succeq0$ (Theorem~\ref{thm:coer-schur}); and the one-sided
asymptotic $\lamA\ge-\varepsilon(T)$ (Proposition~\ref{prop:shape-crit}(e)). A proof of any of them is a proof of \eqref{eq:target} and has
the shape described below.
\end{lemma}

\subsection{The theorem}

\begin{theorem}[the shape of a proof]\label{thm:shape}
Let $\mathcal A$ be a deduction class with property set $\mathsf P$.
\begin{enumerate}
\item[(A)] \textup{(transfer, proved)} If $\mathcal A$ is closed under one of $\mathrm O_1$, $\mathrm O_2$, $\mathrm O_3$, $\mathrm O_4$ (for
all $c>0$), $\mathrm O_6$, $\mathrm O_7$, $\mathrm O_8$, $\mathrm O_{13}$, $\mathrm O_{14}$, $\mathrm O_{15}$, in the ranges of
Lemmas~\ref{lem:shape-del}--\ref{lem:shape-scale}, \ref{prop:shape-eulerdel}--\ref{lem:shape-ft}, \ref{lem:shape-fam}--\ref{lem:shape-poisson},
or if $\mathcal A$ is one of the classes of Lemma~\ref{lem:shape-classes}(i),(ii), then $\mathcal A$ contains no valid deduction of
$\mathrm{WP}_T(\Lambda)$ at the windows named in the corresponding lemma; in particular none proving \eqref{eq:target}.
\item[(A$'$)] \textup{(transfer with a label)} The same conclusion holds for $\mathrm O_4^\varepsilon$ and $\mathrm O_5$ under the
hypotheses $\kappa_T\ne0$, respectively $c_{F_\star}(\log n_0)\ne0$ (conditional); for $\mathrm O_{10}$ for the complete form (cited) and,
for the primitive bottom, argued; for $\mathrm O_{11}$, argued; for $\mathrm O_{12}$ the falsifier is open; and for $\mathrm O_9$ the exclusion
concerns pointwise bounds on the curvature of $\ell$, not the window form.
\item[(B)] \textup{(structure)} Every valid deduction of \eqref{eq:target}, in whatever class, satisfies
Lemmas~\ref{lem:shape-margin}--\ref{lem:shape-restate}: its margin at $T\ge1$ is at most $4\cdot10^9Te^{-15.449T}$; it is exact at the
centre, at the integer resolution and at the density ceiling, and, under the hypothesis of Proposition~\ref{prop:shape-crit}(b), at every
packet frequency; it is not a symmetrization of $D_\sigma$; written in shift-multiplier calculus it is a measure identity on $(-2T,2T)$
plus Bochner positivity, and it is no arithmetic square decomposition for $T\ge\frac12\log3$; it uses the Euler product over $\N$ and the
functional equation at conductor $1$ jointly; and each statement of Lemma~\ref{lem:shape-restate} is \eqref{eq:target} itself.
\item[(C)] \textup{(positive form)} Hence every valid proof of \eqref{eq:target} uses, at a load-bearing line, a property of $\zeta$ that
is destroyed by each modification named in (A), and, with the labels stated there, by those named in (A$'$). Corollary~\ref{cor:shape}
lists these properties.
\end{enumerate}
\end{theorem}
\begin{proof}
(A) is Lemma~\ref{lem:transfer}: each cited lemma exhibits $\mathrm O(q_0)$ with property $\mathsf P$ and a negative bottom at the named
windows, so a deduction in $\mathcal A$ would apply to $\mathrm O(q_0)$ and yield a false conclusion there. (A$'$) is the same argument,
with the falsifier's negativity carrying the label of its lemma. (B) is the conjunction of the structural lemmas. (C) is the
contrapositive of (A): a proof whose property set were invariant under one of the named modifications would lie in a class closed under
it.
\end{proof}

\subsection{Corollaries}

\begin{corollary}[what a proof must be]\label{cor:shape}
A valid proof of \eqref{eq:target}:
\begin{enumerate}
\item[(a)] uses the vector $(\Lambda(n))_{n<e^{2T}}$ exactly at every $T$, not through membership in a class closed under any modification
of (A) (Lemmas~\ref{lem:shape-del}--\ref{lem:shape-scale}, \ref{prop:shape-eulerdel}), and non-linearly: a linear use is a measure identity
on the window (Lemma~\ref{lem:shape-lin}(i)) and so a certificate of class (D), whose content is $\lamP$ itself
(Lemma~\ref{lem:shape-lin}(ii), Lemma~\ref{lem:shape-restate}); and it uses a property not shared by the zero-side modifications
(Lemmas~\ref{lem:shape-quartet}, \ref{lem:shape-ft}): the quartet models keep the comb exactly, so the property is not a property of the
comb alone;
\item[(b)] uses the Euler product over $\N$, the functional equation at conductor $1$ and the integrality of the frequencies jointly
(Proposition~\ref{prop:shape-inputs}; the third input argued), hence is a computation with the data of $\zeta$ at every window, uniform in
$T$, and not a class argument over a family of objects;
\item[(c)] decides a correlation. If $F_\star$ is a real unit minimizer at $T$, then by Lemma~\ref{lem:degen}
$\lambda_{\mathrm A}(T;v')\le\lamA-\sum_n(v'_n-v_n)c_{F_\star}(\log n)$ for every weight vector $v'$: the vector
$(-c_{F_\star}(\log n))_n$ is a supergradient of the concave bottom, and the one-sided derivative in a direction is the minimum of this
expression over the minimizers. No bounded functional of the weights or of the symbol (Lemma~\ref{lem:shape-weight},
Lemma~\ref{lem:shape-margin}(ii), Remark~\ref{rem:wit-sens}), no local criterion (Lemma~\ref{lem:shape-local}) and no termwise bound
(Lemma~\ref{lem:shape-misc}) decides it; numerically, a relative change of $10^{-5}$ in the weights changes the sign at $T=\log2$ while
every coherence functional moves by less (Remark~\ref{rem:wit-sens});
\item[(d)] has no slack: Lemma~\ref{lem:shape-margin}(i) and Proposition~\ref{prop:shape-crit}.
\end{enumerate}
\end{corollary}
\begin{proof}
Each item restates the lemmas it cites; in (c), Lemma~\ref{lem:degen} gives the inequality for every $v'$, which is the definition of a
supergradient of $\mu_T$ at $v_T(\Lambda)$, and a concave function has one-sided directional derivatives equal to the infimum of its
supergradients in that direction.
\end{proof}

\begin{remark}[what the theorem does not say]\label{rem:shape-scope}
(i) It does not say that \eqref{eq:target} is unprovable: every clause excludes a deduction, and the proofs of $\mathrm{WP}_T$ for
$T\le\frac12\log5$ (Theorems~\ref{thm:C}, \ref{thm:log5A}, \ref{thm:log5W}), with the cited ranges of \cite{Zhu2026} and \cite{Liu2026},
are computations with the exact vector, inside the region of Problem~\ref{prob:region} below. (ii) No equivalence between
\eqref{eq:target} and a statement about another object is used; the equivalences of Lemma~\ref{lem:shape-restate} are internal
restatements, used only to say that none of them is weaker than \eqref{eq:target}. (iii) The clauses that are not proved are listed in
Theorem~\ref{thm:shape}(A$'$) and in Proposition~\ref{prop:shape-crit}(b); what would upgrade each is stated with it: a proof that
$\kappa_T\ne0$ or that the degenerate case of Lemma~\ref{lem:degen} does not occur; a primitive witness for $\Xi_Q$; an explicit formula
with compactly supported tests for one Beurling system; a window and a field with a negative primitive bottom; an $m$-fold twist with
a window witness. (iv) The theorem does not rank what remains; it describes the region in which every proof must lie, and that region
is stated as Problem~\ref{prob:region}.
\end{remark}
